\documentclass{article}
\usepackage{tcee-repaso}
\tcetema[Los economistas neoclásicos]{3.A.3}{Los economistas neoclásicos.}
\tcefecha{04/11/2026}
\begin{document}
\cabecerarepaso
\begin{columnas}

% ==========================================================================
\bloque{Esquema del cante}
\esquemacante{%
\apartadocante{Introducción: «nadie sabe lo que tiene\ldots»}{2,5}
\apartadocante{I. Caracterización general de la escuela}{5,5}
\subcante{1. Según \autor{Schumpeter}}{1,5}
\subcante{2. Autores y rango histórico}{2}
\subcante{3. Método}{2}
\apartadocante{II. Revolución marginalista (1871-1874)}{13}
\subcante{1. \autor{Jevons}}{2}
\subcante{2. \autor{Menger} y Viena}{3,5}
\subcante{3. \autor{Walras} y Lausana}{2,5}
\subcante{4. \autor{Marshall} y Cambridge}{5}
\apartadocante{III. Segunda generación}{8}
\subcante{1. Productividad marginal (1890-1899)}{4}
\subcante{2. Interés, dinero y bienestar}{4}
\apartadocante{Conclusión}{1}}
\emph{Reparto}: \autor{Marshall} $\geq$4,5 min (referencia del equilibrio parcial: mercado de pescado en la pizarra); las 3 escuelas, en el método y con cada fundador. \emph{Pizarra}, 3 zonas. \emph{Izda.}: refrán y 5 pegatinas: «lo poco agrada y lo mucho enfada» (UMg decreciente), «cada loco con su tema» (\emph{Methodenstreit}), «una mano lava la otra, y las dos, la cara» (tijeras), «más vale un toma que dos te daré» (impaciencia e interés), «a cada cual lo suyo» (productividad marginal, \autor{Clark}). \emph{Centro}: agua y diamantes, tabla de \autor{Menger}, pescado (se borran por bloque). \emph{Dcha.}: esquema de las 3 escuelas y, recuadradas hasta el final, 2.ª ley de \autor{Gossen} y agotamiento del producto.

% ==========================================================================
\bloque[2,5]{Introducción}
\begin{itemize}
\item \textbf{Enganche}: «Nadie sabe lo que tiene hasta que lo pierde». \autor{Smith} (1776, I, IV): el agua, lo más útil, apenas se cambia por nada; el diamante, casi sin uso, por mucho. Con valor = coste, sin buena respuesta. 1871-1874, 3 autores sin conocerse (Manchester, Viena, Lausana): lo fija la utilidad de la \emph{última} unidad, la \concepto{utilidad marginal}. Pregunta: ¿qué determina el valor: lo que cuesta producir o lo que valoramos?
\item \textbf{Relevancia}: \emph{teoría}: la micro de hoy es su análisis marginal (demanda \vertema{3.A.8}, producción \vertema{3.A.11}, mercados \vertema{3.A.16}, EG \vertema{3.A.21}, bienestar \vertema{3.A.22}); \emph{política}: precios como señales, EC en coste-beneficio, impuesto pigouviano = impuestos al carbono \vertema{3.A.23}, tipo natural de \autor{Wicksell} (lo estiman los bancos centrales) \vertema{3.A.36}; \emph{temario}: la teoría que \autor{Keynes} llama «clásica» \vertema{3.A.4}; 2 de los 4 motivos de \autor{Schumpeter} (1954).
\item \textbf{Contextualización}: paradigmas (\autor{Kuhn}, 1962) \vertema{3.A.1}; ¿revolución? Tardó $\approx$20 años: proceso, no vuelco (\autor{Blaug}, 1972). Clásicos \vertema{3.A.2} (\autor{Smith} 1776 $\to$ \autor{Mill} 1848): valor = coste/trabajo, crecimiento y reparto entre clases; \autor{Mill}: «la teoría está completa». Cabos sueltos: agua-diamantes, bienes no reproducibles (\autor{Ricardo}), valor trabajo $\to$ explotación (\autor{Marx}). 2.ª revolución industrial: PIB pc 1870-1913, R.\,Unido 0,8\,\%, Alemania 1,6\,\%, EE.\,UU. 1,75\,\% (Maddison): de cómo crecer a cómo asignar. Profesión: \autor{Marshall}, cátedra (1885); \emph{Economic Journal} (1891). 1871 \autor{Jevons} (\emph{Theory}) y \autor{Menger} (\emph{Grundsätze}); 1874-77 \autor{Walras} (\emph{Éléments}); 1890 \autor{Marshall} (\emph{Principles}). Centro $\to$ micro; lo agregado, en la teoría monetaria y del ciclo. «Neoclásicos»: \autor{Veblen} (1900), por la continuidad de \autor{Marshall}; \autor{Hicks} y \autor{Stigler} (1941) lo extienden a todos los marginalistas.
\item \textbf{Problemática}: ¿qué cambia frente a los clásicos (valor, objeto, método)? ¿Fue revolución? ¿Qué une y separa a Cambridge, Lausana y Viena? ¿Cómo llega el margen a producción y distribución? ¿Interés, dinero, bienestar; qué queda hoy?
\item \textbf{Estructura}: \hyperref[rep:I]{I}, la escuela (quiénes, cuándo, método, 3 escuelas); \hyperref[rep:II]{II}, los 4 fundadores, de la UMg a las tijeras; \hyperref[rep:III]{III}, distribución (\autor{Wicksteed}, \autor{Wicksell}, \autor{Clark}) e interés, dinero y bienestar (\autor{Fisher}, \autor{Edgeworth}, \autor{Pareto}, \autor{Pigou}).
\end{itemize}

% ==========================================================================
\bloque[5,5]{I. Caracterización general de la escuela\label{rep:I}}
\begin{ideaclave}
Sin un líder: 3 escuelas nacionales (Cambridge, Lausana, Viena) que comparten valor subjetivo, margen y equilibrio, y discrepan en el método: palabras o matemáticas, equilibrio parcial o general, reposo o proceso.
\end{ideaclave}

\sub{I.1. Según \autor{Schumpeter}\label{rep:I.1}}
\begin{itemize}
\item 3 rasgos (\autor{Schumpeter}, 1954): líder, seguidores, cuerpo doctrinal. Doctrina común, pero varios líderes.
\item \emph{Líderes}: \autor{Jevons}, \autor{Menger}, \autor{Walras} y \autor{Marshall}; \autor{Walras} y \autor{Marshall}, padres de la economía neoclásica (juntan oferta y demanda en un sistema).
\item \emph{Seguidores}: \emph{Cambridge} (la neoclásica, no la postkeynesiana \vertema{3.A.4}): \autor{Pigou}; en Inglaterra, \autor{Edgeworth}, \autor{Wicksteed}. \emph{Lausana}: \autor{Pareto}, \autor{Barone}. \emph{Viena}: \autor{Böhm-Bawerk}, \autor{Wieser}; luego \autor{Mises}, \autor{Hayek}. \emph{Fuera}: \autor{Wicksell} (Suecia), \autor{Clark} y \autor{Fisher} (EE.\,UU.).
\item \emph{Objeto}: asignación de recursos escasos entre usos alternativos y precios relativos; distribución = precios de los factores. Se estrecha y se vuelve estático. \autor{Robbins} (1932): relación entre fines y medios escasos con usos alternativos. En lo agregado, ley de \autor{Say}.
\end{itemize}
\begin{center}\footnotesize\setlength{\tabcolsep}{3pt}
\begin{tabular}{@{}>{\raggedright\arraybackslash}p{0.17\linewidth}>{\raggedright\arraybackslash}p{0.34\linewidth}>{\raggedright\arraybackslash}p{0.39\linewidth}@{}}
\toprule
\textbf{Rasgo} & \textbf{Clásicos} & \textbf{Neoclásicos}\\
\midrule
Pregunta & Crecimiento LP, reparto & \textbf{Asignación}, precios relativos\\
Sujeto & Clases sociales & \textbf{Individuo} que optimiza\\
Valor & Objetivo: coste, trabajo & \textbf{Subjetivo}: UMg; luego O y D\\
Distribución & Entre clases; beneficio residuo & \textbf{Entre factores}: PMg\\
Método & Verbal, deductivo, instituciones & \textbf{Margen} y equilibrio\\
Política & \emph{Laissez-faire} & También; \autor{Pigou}: corregir fallos\\
\bottomrule
\end{tabular}
\end{center}

\sub{I.2. Autores y rango histórico\label{rep:I.2}}
\begin{itemize}
\item \emph{De los clásicos}: rendimientos decrecientes (\autor{Turgot}), renta diferencial (\autor{Ricardo}, 1817: margen aplicado a la tierra), deducción, libre competencia.
\item \emph{Precursores} (\concepto{utilidad}: capacidad de satisfacer necesidades): \autor{Bentham} (1789, 1802), \autor{Lloyd} (1833, total vs. marginal); \autor{Bernoulli} (1738): UMg decreciente del dinero, paradoja de San Petersburgo \vertema{3.A.10}; \autor{Cournot} (1838): $D=F(p)$, demanda y oferta en el plano $(q,p)$; el cruce, \autor{Rau} (1841) y \autor{Jenkin} (1870); \autor{Dupuit} (1844): excedente del consumidor de una obra pública \vertema{3.A.16}; von \autor{Thünen} (1850, vol.~2 de \emph{Der isolirte Staat}): pago según la productividad marginal.
\item \autor{Gossen} (1854), 2 leyes: (1) \emph{saturación}: UMg decreciente hasta la saciedad («lo poco agrada y lo mucho enfada»); (2) \emph{igualdad de UMg ponderadas} \vertema{3.A.8}:
\funcionclave{\frac{U_1}{p_1}=\frac{U_2}{p_2}=\dots=\frac{U_n}{p_n}=\lambda}
\item[] La 3.ª, «ley de la escasez», la añaden manuales; \autor{Gossen} no la formula. Pasó inadvertido: \autor{Jevons} le reconoce prioridad (2.ª ed., 1879).
\item \emph{1.ª generación} (\autor{Landreth}-\autor{Colander}): \autor{Jevons}, \autor{Menger}, \autor{Walras}, casi simultáneos e independientes. \emph{5 ideas}: (1) valor \emph{subjetivo}; (2) se decide en el \emph{margen} (escasez); (3) principio \emph{equimarginal} (\autor{Gossen}); (4) utilidad \emph{cardinal} para cada uno, sin comparaciones entre personas; (5) oferta en segundo plano: el valor de los factores se deriva del de los bienes.
\item \autor{Jevons}, \autor{Menger}: cadenas causales; \autor{Walras}: interdependencia. Juntan las 2 hojas \autor{Marshall} (un mercado) y \autor{Walras} (todos): fin del valor trabajo y de la distribución residual.
\item \emph{3 etapas} (\autor{Fonseca}): (1) revolución 1871-1874 (+ \autor{Marshall}, 1890); (2) distribución por productividad marginal, 1890-1894 (\autor{Clark}, \autor{Wicksteed}, \autor{Wicksell}; \autor{Clark} 1899); (3) \emph{resurgimiento paretiano}, 1934-1947 (\autor{Hicks}-\autor{Allen}, \autor{Hotelling}, \autor{Lange}, \autor{Allais}, \autor{Samuelson}). Entre 2.ª y 3.ª: \autor{Fisher}, \autor{Edgeworth}, \autor{Pareto}, \autor{Pigou}.
\end{itemize}

\sub{I.3. Método\label{rep:I.3}}
\begin{itemize}
\item \emph{Común}: \concepto{individualismo metodológico}; \concepto{análisis marginal} (optimización); \concepto{equilibrio} (nadie quiere cambiar); \concepto{abstracción}: hogares que maximizan utilidad, empresas beneficio; deducción y cálculo (en el texto en Lausana, en notas en Cambridge, sin él en Viena).
\item \emph{Lo que separa}, 3 preguntas: ¿palabras o matemáticas? ¿Un mercado o todos? ¿Reposo o proceso?
\end{itemize}
\begin{center}\footnotesize\setlength{\tabcolsep}{3pt}
\begin{tabular}{@{}>{\raggedright\arraybackslash}p{0.15\linewidth}>{\raggedright\arraybackslash}p{0.26\linewidth}>{\raggedright\arraybackslash}p{0.26\linewidth}>{\raggedright\arraybackslash}p{0.26\linewidth}@{}}
\toprule
 & \textbf{Cambridge} & \textbf{Lausana} & \textbf{Viena}\\
\midrule
Fundador & \autor{Marshall} 1890 & \autor{Walras} 1874 & \autor{Menger} 1871\\
Método & \textbf{Verbal}; mates en notas & \textbf{Matemático} & \textbf{Verbal}, deductivo\\
Equilibrio & \textbf{Parcial}, plazos & \textbf{General} & \textbf{Proceso}\\
Valor & Utilidad y coste: tijeras & \emph{Rareté}, O y D de todo & Subjetivo; imputación\\
Interés & Cuasirrentas; espera & Capital en el sistema & Rodeos; preferencia temporal\\
Dinero & $M^D=kPY$ & Circulación (última etapa) & Origen; ciclo\\
Bienestar & EC; \autor{Pigou} & Óptimo de \autor{Pareto} & Crítica del cálculo socialista\\
\bottomrule
\end{tabular}
\end{center}
\begin{itemize}
\item \concepto{Disputa sobre el método} (\emph{Methodenstreit}) \vertema{3.A.1}: \autor{Menger}, \emph{Untersuchungen} (1883), teoría abstracta vs. escuela histórica alemana; \autor{Schmoller} (1883), reseña durísima (economía inductiva, historia e instituciones); \autor{Menger}, \emph{Die Irrthümer des Historismus} (1884). «Cada loco con su tema». Austriacos: deductivos y sin matemáticas, distintos de historicistas y de \autor{Walras}.
\end{itemize}

% ==========================================================================
\bloque[13]{II. Revolución marginalista (1871-1874)\label{rep:II}}
\begin{ideaclave}
El valor pasa a depender de la utilidad marginal; \autor{Menger} lo lleva a los factores y al tiempo, \autor{Walras} a todos los mercados a la vez y, en 1890, \autor{Marshall} une utilidad y coste en las tijeras de la oferta y la demanda.
\end{ideaclave}

\sub{II.1. \autor{William Stanley Jevons} (1835-1882)\label{rep:II.1}}
\begin{itemize}
\item \emph{Contexto}: químico y ensayador (casa de la moneda de Sídney), lógico, estadístico; Manchester y Londres. Famoso por \emph{The Coal Question} (1865). \emph{The Theory of Political Economy} (1871). \emph{Método}: matemático y utilitarista (cálculo del placer y el dolor, \autor{Bentham}).
\item \emph{Utilidad aditiva}; cuenta el \concepto{grado final de utilidad} (UMg), decreciente:
\begin{formula}
\centering $U(\vect x)=u_1(x_1)+\dots+u_n(x_n),\quad U_i=u_i'(x_i),\ u_i''<0$
\end{formula}
\item[] UMg de un bien no depende de los demás: sin sustitutivos ni complementarios. Reparte la renta según \autor{Gossen}.
\item \emph{Agua y diamantes}, resuelta: mucha más utilidad total del agua, pero tanta que la última vale poco; la marginal fija el precio.
\end{itemize}
\graficorepaso[0.85\linewidth]{agua-diamantes}{Área bajo la curva hasta $\bar x$ = utilidad total; altura en $\bar x$ = UMg, la que fija el precio}
\begin{itemize}
\item \emph{Ley de la indiferencia}: en un mercado abierto, en un momento, no puede haber 2 precios para el mismo artículo: \concepto{ley del precio único}.
\item \emph{Intercambio} ($A$ con bien 1, $B$ con bien 2):
\begin{formula}
\centering $\mathrm{RMS}^A_{12}=\dfrac{U^A_1}{U^A_2}=\dfrac{x_2}{x_1}=\dfrac{U^B_1}{U^B_2}=\mathrm{RMS}^B_{12}$
\end{formula}
\item[] Con 2 personas, indeterminado (\autor{Edgeworth}).
\item \emph{Cadena de \autor{Jevons}} (1879): coste $\to$ oferta $\to$ grado final de utilidad $\to$ valor. El coste importa solo por la cantidad.
\item \emph{Trabajo}: hasta que la penosidad de la última hora = utilidad de lo obtenido \vertema{3.A.25}.
\item \emph{Otras}: salarios y beneficios no van en sentido contrario (2 incógnitas con 1 ecuación; fondo de salarios «ilusorio»); \concepto{paradoja de \autor{Jevons}} (1865): más eficiencia en el carbón ↑ su consumo (efecto rebote); \emph{ciclo} por manchas solares (1875, 1878): teoría \emph{real}; ojo: los \emph{sunspot equilibria} (\autor{Azariadis}, 1981; \autor{Cass}-\autor{Shell}, 1983) no son suyos \vertema{3.A.42}.
\item \emph{Valoración}: la UMg más clara en matemáticas; aditividad y desinterés por la oferta: no explica el coste.
\end{itemize}

\sub{II.2. \autor{Carl Menger} (1840-1921) y Viena\label{rep:II.2}}
\begin{itemize}
\item \emph{Contexto}: periodista y funcionario; \emph{Grundsätze} (1871); Univ. de Viena (1873): escuela austriaca. \emph{Método}: deductivo, verbal, sin matemáticas (la «praxeología» es de \autor{Mises}).
\item \emph{Bienes}: \concepto{de orden inferior} (pan) y \concepto{de orden superior} (harina, horno, trabajo), complementarios. Su valor se \concepto{imputa} desde el de los de consumo: no fijan los costes el precio, sino al revés.
\item \emph{Valor de una unidad}: importancia de la necesidad menos urgente que depende de ella (UMg sin la palabra). Escasez relativa al deseo: huevos podridos (\autor{Robbins}).
\end{itemize}
\begin{center}\footnotesize
\begin{tabular}{@{}c*{5}{c}@{}}
\toprule
\textbf{Unidad} & \textbf{I} & \textbf{II} & \textbf{III} & \textbf{IV} & \textbf{V}\\
\midrule
1.ª & \textbf{10} & \textbf{9} & \textbf{8} & 7 & 6\\
2.ª & \textbf{9} & \textbf{8} & 7 & 6 & 5\\
3.ª & \textbf{8} & 7 & 6 & 5 & 4\\
4.ª & 7 & 6 & 5 & 4 & 3\\
\bottomrule
\end{tabular}\\[1pt]
{\color{tceegris}\scriptsize\itshape Escala de \autor{Menger}: con 6 unidades, 3 a I, 2 a II, 1 a III; la última cubre 8 en cada clase: valor 8}
\end{center}
\begin{itemize}
\item \emph{Origen del dinero} (1892): no lo inventó el Estado; se acepta el bien más vendible, que acaba siendo medio de cambio de todos: institución no diseñada \vertema{3.A.35}.
\item \textbf{\autor{Böhm-Bawerk}} (1851-1914), \emph{Kapital und Kapitalzins} (1884, 1889): rechaza productividad, abstinencia (\autor{Senior}) y explotación (\autor{Marx}); interés = \emph{agio}: los bienes presentes valen más. «Más vale un toma que dos te daré». 3 causas: (1) mejor provisto en el futuro; (2) se infravalora el futuro; (3) superioridad técnica: \concepto{rodeos de producción} (red vs. mano), mayor \emph{periodo medio de producción}. 1 y 2, preferencia temporal; 3, técnica (\autor{Fisher} la rechaza). Periodo medio no medible con muchos bienes de capital: anticipo de la controversia del capital \vertema{3.A.4}. \emph{Zum Abschluss des Marxschen Systems} (1896): valor del libro~I de \emph{El capital} incompatible con los precios de producción del libro~III \vertema{3.A.2}.
\item \textbf{\autor{Wieser}} (1851-1926): acuña \emph{Grenznutzen} (1884); imputación como teoría general (1889); \emph{coste alternativo} como base del coste; \emph{opportunity cost}: \autor{Green} (1894).
\item \textbf{\autor{Mises}} (1881-1973): \emph{Theorie des Geldes} (1912): UMg al dinero (poder adquisitivo de ayer, hasta la mercancía). \emph{Ciclo austriaco} (\autor{Mises}, 1912; \autor{Hayek}, 1931): crédito bajo el tipo natural $\Rightarrow$ procesos más largos que el ahorro voluntario; al frenarse, crisis que liquida errores: evitar el auge \vertema{3.A.42}. \emph{Cálculo socialista} (1920): sin mercados de capital, sin precios; \autor{Lange} (1936): el planificador tantea como \autor{Walras}.
\item \textbf{\autor{Hayek}} (1899-1992): conocimiento disperso que resumen los precios (1945); competencia como proceso de descubrimiento (1948) \vertema{3.A.16}. Línea de \autor{Schumpeter} (1912) y \autor{Kirzner} (1973). Nobel 1974.
\item \emph{Valoración}: el subjetivismo más coherente (tiempo, capital, dinero, conocimiento), sin matemáticas ni sistema cerrado.
\end{itemize}

\sub{II.3. \autor{Léon Walras} (1834-1910) y Lausana\label{rep:II.3}}
\begin{itemize}
\item \emph{Contexto}: hijo de \autor{Auguste Walras} (\emph{rareté}); \autor{Cournot}; Lausana desde 1870; \emph{Éléments} (1874, 1877); socialista. \emph{Método}: economía \emph{pura} (valor de cambio, como la mecánica, con matemáticas), \emph{aplicada} (producción) y \emph{social} (justicia).
\item \emph{Idea}: todo depende de todo: ¿precios que vacían todos los mercados a la vez? \vertema{3.A.21} \emph{Supuestos}: competencia; consumidores maximizan y venden servicios de factores; coeficientes de fabricación; \concepto{numerario}. 4 etapas: intercambio, producción, capitalización, circulación (dinero).
\item \emph{Rareté} = UMg; proporcionales a los precios (\autor{Gossen}) $\Rightarrow$ demandas en función de todos los precios (versión general de la demanda recíproca de \autor{Mill}).
\item \concepto{Ley de \autor{Walras}}: por las restricciones presupuestarias, valor de los excesos de demanda = 0 a cualquier precio; si $n-1$ mercados en equilibrio, el último también; solo precios relativos.
\item Cuenta ecuaciones e incógnitas: iguales $\Rightarrow$ determinado. \concepto{Tanteo} (\emph{tâtonnement}): precio ↑ con exceso de demanda, ↓ con exceso de oferta. No habla de subastador («on y crie»); en la 4.ª ed. (1900), \emph{bons} tras las críticas de \autor{Bertrand} (1883) y \autor{Edgeworth}.
\item \emph{Implicaciones}: precios y cantidades a la vez, sin cadena causal; empresario sin ganancia ni pérdida; precio del capital = valor capitalizado de sus servicios.
\item \emph{Valoración}: contar ecuaciones no basta (soluciones no negativas): existencia, \autor{Wald} (1936) y \autor{Arrow}-\autor{Debreu} (1954); convergencia del tanteo no garantizada. Traducción de \autor{Jaffé} (1954). Marco de la micro y de buena parte de la macro.
\item \textbf{\autor{Pareto}} (1848-1923), sucede a \autor{Walras} (1893); \emph{Cours} (1896), \emph{Manuale} (1906): \emph{ofelimidad}, utilidad ordinal \vertema{3.A.8}. \textbf{\autor{Barone}} (1908): un ministerio de la producción llegaría al resultado de la competencia: origen del debate del cálculo socialista.
\end{itemize}

\sub{II.4. \autor{Alfred Marshall} (1842-1924) y Cambridge\label{rep:II.4}}
\begin{itemize}
\item \emph{Contexto}: matemático, llega desde la ética y la pobreza; catedrático 1885-1908; \emph{Principles of Economics} (1890, 8 ed. hasta 1920) sustituyen a \autor{Mill}.
\item \emph{Método}: «estudio de la humanidad en los asuntos corrientes de la vida»; \emph{equilibrio parcial} (corral del \emph{ceteris paribus}) \vertema{3.A.16}; mates en notas (carta a \autor{Bowley}, 1906: traducir a palabras y «quemar las matemáticas»); continuidad: \emph{natura non facit saltum}.
\item \emph{Frente a \autor{Walras}}: la cantidad es la variable independiente; \emph{precio de demanda} $p^D(q)$ $\Rightarrow$ precio en el eje vertical \vertema{3.A.9}.
\item \emph{Idea}: las \emph{tijeras}: discutir si corta la hoja de arriba o la de abajo = si rige la utilidad o el coste. «Una mano lava la otra, y las dos, la cara». Plazo corto: pesa la demanda; largo: el coste.
\item \emph{Supuestos}: mercado pequeño, competencia, \emph{UMg del dinero constante} (hoy, utilidad cuasilineal); \concepto{empresa representativa}.
\item \emph{Demanda}: precio de demanda = UMg en dinero, decreciente (ley de la demanda). \concepto{Elasticidad} $\varepsilon_D=\frac{dq}{dp}\frac pq$; \concepto{excedente del consumidor}:
\begin{formula}
\centering $\mathrm{EC}=\displaystyle\int_0^{q}p^D(s)\,ds-p\,q$
\end{formula}
\item \emph{Oferta}: $p^S(q)$ cubre el coste; economías \emph{internas} (tamaño de la empresa) y \emph{externas} (de la industria; distritos industriales).
\item \emph{Periodos}: \concepto{de mercado} ($q$ dada, precio por la demanda); \concepto{corto plazo} (equipo dado; \concepto{cuasirrenta}: lo que obtiene el equipo fijo sobre el coste variable); \concepto{largo plazo} (equipo y n.º de empresas; \concepto{valor normal} = coste de la empresa representativa); \emph{secular} (técnica, población, capital).
\end{itemize}
\graficorepaso[0.62\linewidth]{marshall-periodos}{Mercado de pescado: $D_0\to D_1$ (epidemia del ganado); $MS$ vertical $\Rightarrow$ equilibrio temporal $B$ (cuasirrentas); salen barcos $\Rightarrow$ equilibrio normal $C$, con $p_A<p_C<p_B$}
\begin{itemize}
\item Con rendimientos crecientes en la pesca, oferta normal decreciente (industria de costes decrecientes) \vertema{3.A.16}.
\item \emph{Ajuste marshalliano}: $q$↑ si precio de demanda $>$ precio de oferta (por cantidades) vs. \emph{walrasiano} (por precios). En el día manda el precio; de un día a otro, la cantidad. En \autor{Marshall}, precios antes que cantidades; \autor{Keynes} invierte el orden (\autor{Leijonhufvud}, 1968) \vertema{3.A.4}.
\item \emph{Valoración}: base de la micro de mercados. \autor{Sraffa} (1926): rendimientos crecientes internos incompatibles con la competencia; externos, con el equilibrio parcial $\to$ competencia imperfecta \vertema{3.A.18}.
\item \emph{Distribución}: O y D de cada factor; 4.º factor, la \emph{organización}. Demanda: \emph{producto neto marginal}; oferta: coste (formar al trabajador; la \emph{espera}, en vez de la abstinencia de \autor{Senior}). \emph{Renta}: factor de oferta fija; \emph{cuasirrenta}: equipo fijo a CP \vertema{3.A.25}.
\item \emph{Dinero}: \concepto{saldos de caja de Cambridge}, $M^D=kPY$, $k=1/V$ (\autor{Marshall}, 1923; formaliza \autor{Pigou}, 1917) \vertema{3.A.35}.
\item \emph{Comercio}: \emph{curvas de oferta recíproca} (\emph{Pure Theory of Foreign Trade}, 1879): relación real de intercambio \vertema{3.B.5}; arancel óptimo \vertema{3.B.7}.
\item \textbf{\autor{Pigou}} (1877-1959), sucesor en 1908: formaliza los saldos de caja; excedente $\to$ bienestar; \emph{Theory of Unemployment} (1933): paro se cura bajando salarios, blanco de \autor{Keynes} \vertema{3.A.4}.
\end{itemize}

% ==========================================================================
\bloque[8]{III. 2.ª generación: distribución, dinero y bienestar\label{rep:III}}
\begin{ideaclave}
En los años 90 el margen llega a la oferta: cada factor cobra su productividad marginal y, con rendimientos constantes, el producto se reparte entero. Sobre esa base, las teorías del interés, del dinero y del bienestar.
\end{ideaclave}

\sub{III.1. Productividad marginal (1890-1899)\label{rep:III.1}}
\begin{itemize}
\item Fechas (\autor{Fonseca}): 1890-1894: \autor{Clark} (1891), \autor{Wicksell} (1893), \autor{Wicksteed} (1894); \emph{The Distribution of Wealth} (\autor{Clark}), 1899.
\item \emph{Piezas}: \concepto{función de producción} $q=F(\vect z)$, PMg positivas y decrecientes \vertema{3.A.11}; competencia perfecta: $w_i=p\,\mathrm{PMg}_i$ (demanda de factor = VPMg) \vertema{3.A.25}; capital = un factor más; interés separado del beneficio, nulo en equilibrio. Sin renta residual, ¿suman el producto?
\item \textbf{\autor{Clark}} (1847-1938): formado con \autor{Knies}; marginalismo en EE.\,UU.; fundador de la AEA; generaliza la renta diferencial a todos los factores (1891). \emph{Ley de la productividad final}: salario = productividad del último trabajador; «no supone explotación»; «a cada cual lo suyo». Estática: beneficio nulo; beneficio positivo, fenómeno dinámico. \emph{Crítica}: la PMg mide lo que se pierde con una unidad menos, no lo que «crea»; depende de los demás factores y del reparto de la propiedad.
\item \textbf{Agotamiento del producto}: \autor{Clark} no lo demostró; \autor{Wicksteed} (1894), con rendimientos constantes y competencia perfecta; \concepto{teorema de \autor{Euler}} $\sum_i\mathrm{PMg}_iz_i=F(\vect z)$; con $w_i=p\,\mathrm{PMg}_i$:
\funcionclave{\sum_i w_i\,z_i=p\sum_i\mathrm{PMg}_i\,z_i=p\,F(\vect z)}
\item[] Sin residuo, beneficio extraordinario nulo: la distribución, un problema de precios. \autor{Flux} (1894) ve que es \autor{Euler}. Rendimientos crecientes: costaría más de lo producido; decrecientes: residuo.
\item \textbf{\autor{Wicksell}} (1851-1926): \autor{Böhm-Bawerk} + \autor{Walras}; tramos crecientes, constantes y decrecientes: con libre entrada, mínimo del CMe (rendimientos localmente constantes) $\Rightarrow$ se agota igual \vertema{3.A.12}.
\item \emph{Valoración}: \autor{Pareto} y \autor{Edgeworth} critican la dependencia de rendimientos constantes; \autor{Wicksteed} la retira. Fondo: el capital solo se suma en valor, que depende del tipo de interés: \emph{controversia del capital}, \emph{retorno de técnicas} \vertema{3.A.4}. \autor{Wicksteed} (1910): coste = coste de oportunidad; relaciones \emph{no tuistas}.
\end{itemize}

\sub{III.2. Interés, dinero y bienestar\label{rep:III.2}}
\begin{itemize}
\item \textbf{\autor{Fisher}} (1867-1947), Yale: \emph{Mathematical Investigations} (1892); interés: 1896, 1907, \emph{Theory of Interest} (1930); dinero: \emph{Purchasing Power of Money} (1911), \emph{Index Numbers} (1922); deflación por deudas (1933).
\item \emph{Interés}: acepta 2 causas de \autor{Böhm-Bawerk} (la \concepto{impaciencia}), rechaza la técnica; impaciencia + oportunidad de invertir (dedica 1930 a \autor{Rae} y \autor{Böhm-Bawerk}):
\begin{formula}
\centering $u'(C_1)=\beta(1+r)\,u'(C_2),\qquad F'(K_2)-\delta=r$
\end{formula}
\item[] \vertema{3.A.33}. \emph{Teorema de separación}: con mercado de capitales perfecto, la inversión no depende de los gustos (máx. valor actual, luego consumo) \vertema{3.A.29}; anticipa a \autor{Modigliani}-\autor{Miller} (1958).
\item \emph{Teoría cuantitativa} (1911), $T$ = transacciones: identidad, teoría si $V$ estable y $T$ real \vertema{3.A.35}:
\begin{formula}
\centering $MV\equiv PT;\qquad 1+i=(1+r)(1+\pi^e)\Rightarrow i\simeq r+\pi^e$
\end{formula}
\item[] Con $k=1/V$, = Cambridge (en \autor{Fisher}, $V$ por hábitos de pago; en Cambridge, $k$ es decisión). Ecuación de \autor{Fisher}, ya en 1896. \emph{Índice ideal}: media geométrica de \autor{Laspeyres} y \autor{Paasche} \vertema{3.A.9}. \emph{Deflación por deudas}: «cuanto más pagan los deudores, más deben» \vertema{3.A.42}.
\item \textbf{\autor{Wicksell}}, \emph{Geldzins und Güterpreise} (1898): el dinero llega a los precios por el tipo. \concepto{Tipo natural} $r^n$: neutral respecto de los precios, iguala ahorro e inversión. $i<r^n$: $I>S$, crédito y precios ↑ mientras dure: \concepto{proceso acumulativo}; $i>r^n$: ↓. Regla: ↑ tipo si suben los precios: antecedente de las reglas de tipos \vertema{3.A.36} y de la NEK \vertema{3.A.7}. Escuela de Estocolmo (\autor{Myrdal}, \autor{Lindahl}, \autor{Lundberg}, \autor{Ohlin}); \emph{unanimidad} (1896) \vertema{3.A.24} $\to$ precios de \autor{Lindahl} (1919) \vertema{3.A.23}.
\item \textbf{\autor{Edgeworth}} (1845-1926), Oxford, 1.er director del \emph{EJ}; \emph{Mathematical Psychics} (1881): utilidad no aditiva $\Rightarrow$ primeras \concepto{curvas de indiferencia} \vertema{3.A.8}; \concepto{curva de contrato}: precio indeterminado en el intercambio bilateral; conjetura que desaparece con muchos agentes: \emph{núcleo} (\autor{Debreu}-\autor{Scarf}, 1963) \vertema{3.A.21}. La caja la dibujó \autor{Pareto} (1906).
\item \textbf{Óptimo de \autor{Pareto}}: no se puede mejorar a nadie sin empeorar a otro; sin comparaciones interpersonales; la competencia lleva a él (1906): germen del 1TFEB \vertema{3.A.22}.
\item \textbf{\autor{Pigou}}: \emph{Wealth and Welfare} (1912) $\to$ \emph{The Economics of Welfare} (1920). \emph{Bienestar económico}: medible con dinero; ↑ con el dividendo nacional y con mejor reparto (UMg decreciente). Producto neto marginal privado $\neq$ social: \concepto{externalidad}; se produce de más:
\begin{formula}
\centering $\mathrm{CMgS}=\mathrm{CMgP}+\mathrm{CMgE},\qquad t^*=\mathrm{CMgE}(q^*)$
\end{formula}
\item[] \concepto{Impuesto pigouviano} \vertema{3.A.23}. \emph{Críticas}: \autor{Robbins} (1932): comparaciones interpersonales $\to$ nueva economía del bienestar (solo \autor{Pareto}); \autor{Coase} (1960): daño recíproco y negociable.
\item \emph{Las piezas juntas}: con precios y salarios flexibles, los mercados se vacían; paro voluntario o por rigideces; dinero solo fija $P$: el modelo «clásico» de \autor{Keynes} (\autor{Heijdra}, 2017) \vertema{3.A.4}.
\item \emph{Valoración}: la 2.ª generación cierra el sistema. Críticas desde fuera: institucionalistas (\autor{Veblen}, 1898: calculadora de placeres y dolores), \autor{Sraffa} (1926), \autor{Keynes} (1936); desde dentro, austriacos (equilibrio sin tiempo ni conocimiento).
\end{itemize}

% ==========================================================================
\bloque[1]{Conclusión}
\begin{itemize}
\item \textbf{Recapitulación}: «En conclusión,» 3 ideas: valor por UMg, no trabajo (1871-1874); O y D a la vez, en un mercado (\autor{Marshall}) o en todos (\autor{Walras}); con la productividad marginal, el margen llega a la distribución, y sobre él interés, dinero y bienestar.
\item \textbf{Relevancia}: margen, optimización y equilibrio, instrumental de la micro; tipo natural e impuesto pigouviano, en la política.
\item \textbf{Extensiones}: resurgimiento paretiano (\autor{Hicks}-\autor{Allen}, 1934; \autor{Hicks}, 1939; \autor{Samuelson}, 1947) \vertema{3.A.8}; \autor{Arrow}-\autor{Debreu} (1954) \vertema{3.A.21}; competencia imperfecta (1933) \vertema{3.A.18}; \autor{Keynes} y los «clásicos» \vertema{3.A.4}; racionalidad limitada y conductual \vertema{3.A.1}.
\item \textbf{Opinión}: siguen vivas la crítica de Cambridge a la PMg del capital agregado y la austriaca al equilibrio estático (\autor{Hayek}, 1948); el instrumental ha sobrevivido a sus críticos.
\item \textbf{Idea final}: ¿cuesta o se valora? Las 2, como las hojas de unas tijeras, y en el margen. \autor{Mill} la creía cerrada en 1848; 23 años después, el agua y los diamantes la reabrieron. Los clásicos fundan la ciencia; los neoclásicos la consolidan.
\end{itemize}

% ==========================================================================
\bloque{Autores y fechas}
\hito{1738}{\autor{Bernoulli}: paradoja de San Petersburgo}
\hito{1776}{\autor{Smith}: paradoja del agua y los diamantes}
\hito{1817}{\autor{Ricardo}: renta diferencial}
\hito{1838}{\autor{Cournot}: $D=F(p)$}
\hito{1844}{\autor{Dupuit}: excedente del consumidor}
\hito{1848}{\autor{Mill}: «la teoría está completa»}
\hito{1850}{von \autor{Thünen}: productividad marginal (vol.~2)}
\hito{1854}{\autor{Gossen}: 2 leyes}
\hito{1865}{\autor{Jevons}: \emph{The Coal Question}}
\hito{1871}{\autor{Jevons}: \emph{Theory}; \autor{Menger}: \emph{Grundsätze}}
\hito{1874}{\autor{Walras}: \emph{Éléments} (2.ª entrega, 1877)}
\hito{1879}{\autor{Marshall}: \emph{Pure Theory of Foreign Trade}}
\hito{1881}{\autor{Edgeworth}: \emph{Mathematical Psychics}}
\hito{1883}{\autor{Menger}: \emph{Untersuchungen}; \autor{Schmoller}: reseña}
\hito{1884}{\autor{Wieser}: \emph{Grenznutzen}; \autor{Böhm-Bawerk}: vol.~1}
\hito{1889}{\autor{Böhm-Bawerk}: \emph{Positive Theorie}; \autor{Wieser}}
\hito{1890}{\autor{Marshall}: \emph{Principles}}
\hito{1892}{\autor{Menger}: origen del dinero}
\hito{1894}{\autor{Wicksteed}: agotamiento; \autor{Flux}: \autor{Euler}}
\hito{1896}{\autor{Fisher}: $i\simeq r+\pi^e$; \autor{Pareto}: \emph{Cours}}
\hito{1898}{\autor{Wicksell}: \emph{Geldzins und Güterpreise}}
\hito{1899}{\autor{Clark}: \emph{The Distribution of Wealth}}
\hito{1900}{\autor{Veblen}: «neoclásicos»}
\hito{1906}{\autor{Pareto}: \emph{Manuale}; óptimo}
\hito{1908}{\autor{Barone}: ministerio de la producción}
\hito{1911}{\autor{Fisher}: $MV\equiv PT$}
\hito{1912}{\autor{Mises}: \emph{Theorie des Geldes}; \autor{Pigou}: \emph{Wealth and Welfare}}
\hito{1917}{\autor{Pigou}: saldos de caja}
\hito{1920}{\autor{Pigou}: \emph{Economics of Welfare}; \autor{Mises}: cálculo}
\hito{1926}{\autor{Sraffa}: crítica a \autor{Marshall}}
\hito{1930}{\autor{Fisher}: \emph{The Theory of Interest}}
\hito{1931}{\autor{Hayek}: \emph{Prices and Production}}
\hito{1932}{\autor{Robbins}: definición; crítica al bienestar}
\hito{1933}{\autor{Fisher}: deflación por deudas}
\hito{1936}{\autor{Wald}: existencia; \autor{Lange}: socialismo}
\hito{1945}{\autor{Hayek}: el uso del conocimiento}
\hito{1954}{\autor{Arrow}-\autor{Debreu}; trad. de \autor{Jaffé}}

% ==========================================================================
\bloque{Datos clave}
\begin{itemize}
\item PIB por habitante, 1870-1913, \% anual: R.\,Unido \dato{0,8}{Maddison}; Alemania \dato{1,6}{Maddison}; EE.\,UU. \dato{1,75}{Maddison}, que supera al R.\,Unido desde 1880.
\end{itemize}

% ==========================================================================
\bloque{Preguntas del tribunal}
\begin{enumerate}
\item \textbf{¿Por qué «neoclásicos»?} \autor{Veblen} (1900), por la continuidad de \autor{Marshall} con los clásicos; \autor{Hicks} y \autor{Stigler} lo extienden a todos; \autor{Keynes} los llama «clásicos».
\item \textbf{¿Fue una revolución?} En teoría del valor, objeto y método, sí; en el tiempo, un proceso de $\approx$20 años (\autor{Blaug}, 1972).
\item \textbf{¿Ajuste marshalliano vs. walrasiano?} Por cantidades (precio de demanda vs. de oferta) / por precios (excesos de demanda).
\item \textbf{¿Se agota siempre el producto?} Solo con rendimientos constantes (\autor{Euler}); \autor{Wicksell}: también en el mínimo del CMe con libre entrada.
\item \textbf{¿Qué causas del interés acepta \autor{Fisher}?} Las 2 subjetivas de \autor{Böhm-Bawerk} (impaciencia); rechaza la técnica (rodeos).
\item \textbf{¿Ley de \autor{Walras}?} Valor de los excesos de demanda = 0 a cualquier precio: $n-1$ mercados en equilibrio $\Rightarrow$ el $n$-ésimo; solo precios relativos.
\end{enumerate}

\begin{trampa}
\begin{itemize}
\item Neoclásicos: utilidad \emph{cardinal} para cada uno, pero \emph{sin} comparaciones entre personas (\autor{Menger} ni la mide). Ordinal: \autor{Pareto} (ofelimidad), \autor{Hicks}-\autor{Allen} (1934).
\item \autor{Cournot} (1838) escribe $D=F(p)$ y no funda la demanda en la utilidad; el cruce de O y D: \autor{Rau} (1841), \autor{Jenkin} (1870).
\item La caja de \autor{Edgeworth} la dibujó \autor{Pareto} (1906); \autor{Edgeworth}: curvas de indiferencia y de contrato.
\item \autor{Walras} no habla de subastador; los \emph{Éléments}: 1874 \emph{y} 1877.
\item «Praxeología»: \autor{Mises}, no \autor{Menger}. \emph{Grenznutzen}: \autor{Wieser}; \emph{opportunity cost}: \autor{Green} (1894).
\item 3.ª etapa de \autor{Fonseca} (1934-1947): \autor{Hicks}-\autor{Allen}, \autor{Samuelson}\ldots, no \autor{Fisher}, \autor{Edgeworth}, \autor{Pareto}, \autor{Pigou}.
\item $MV\equiv PT$: identidad; $T$ = transacciones, no impuestos. $k=1/V$.
\item Productividad marginal: \autor{Thünen} (1850) antes que \autor{Clark}; agotamiento demostrado por \autor{Wicksteed}, no \autor{Clark}.
\end{itemize}
\end{trampa}

\end{columnas}
\end{document}
